\section{Purpose of the Project}
\label{sec:purpose}

\subsection{Goals}
\label{sec:goals}

The project generates the Kubernetes files of a small cluster from short
deployment documents. Each application's repository holds one \emph{project
document}, which states what the application needs, and one \emph{platform
document} states what all applications share. Task~1 defined two metamodels in
Ecore: the \emph{source metamodel}, which describes both documents, and the
\emph{target metamodel}, which describes the \emph{resolved model}. A resolved
model records every value the platform derives for one project, and the files
that will be generated from it. Task~2 implemented the model transformation
from source models to a resolved model in QVT-Operational.

Task~3 implements the last step: Acceleo templates that turn a resolved model
into text. The text is the \emph{Deliverable Set}, the YAML and JSON files a
cluster applies: the workloads and their service accounts, the network
policies, the routes of the web proxy, the monitors, the volume claims and
backups, the configuration files, the objects of the secret operator, and the
policies of the secret store. The templates are run on the resolved models the
Task~2 transformation writes, so the three tasks together form one pipeline
from documents to files.

The code generation has two goals.

\begin{enumerate}
  \item It writes the same files as the production implementation of the same
    specification, which is written in TypeScript and used to deploy the
    cluster, \emph{byte for byte}. Both implementations are compared with the
    same committed trees of files, called \emph{oracles}. Byte-for-byte
    equality is what makes the comparison meaningful: whitespace, key order and
    quoting decide whether two files are the same, so the templates have to
    decide them exactly as the production implementation's serializer does.
  \item It decides nothing. Every value in a generated file is read from the
    resolved model. A template only chooses how a value is spelled in YAML or
    JSON, never which value it is, so a wrong file is always traced to the
    transformation or the model and never to the templates.
\end{enumerate}

\subsection{Role of the Code Generation}
\label{sec:role}

Figure~\ref{fig:pipeline} shows the pipeline. Xtext parses each document into a
source model and OCL rejects invalid ones (Task~1). Two QVT-Operational
transformations lower and resolve them into one resolved model per project
(Task~2). The Acceleo module of this report writes the files of every resolved
model it is given. One Java entry point, \texttt{Pipeline.render}, runs all of
these steps in one pass, and the build compares the files it writes with the
oracles.

\begin{figure}[!htbp]
\centering
\begin{tikzpicture}[
  node distance=0.45cm and 0.55cm,
  box/.style={draw, rounded corners=2pt, align=center, font=\scriptsize,
    minimum height=0.9cm, text width=2.1cm},
  document/.style={box, fill=blue!6},
  stage/.style={box, fill=green!8},
  result/.style={box, fill=orange!10},
  arr/.style={-{Stealth[length=2mm]}}]
\node[document] (docs) {Project and platform documents (YAML)};
\node[stage, right=of docs] (parse) {Xtext parsing and OCL (Task~1)};
\node[stage, right=of parse] (qvto) {Lowering and resolution (QVTo, Task~2)};
\node[result, right=of qvto] (model) {Resolved models (XMI)};
\node[stage, right=of model] (acceleo) {\textbf{Code generation (Acceleo)}};
\node[result, below=of acceleo] (tree) {The Deliverable Set (YAML, JSON)};
\node[result, below=of model] (oracle) {Committed trees, compared byte for byte};
\draw[arr] (docs) -- (parse);
\draw[arr] (parse) -- (qvto);
\draw[arr] (qvto) -- (model);
\draw[arr] (model) -- (acceleo);
\draw[arr] (acceleo) -- (tree);
\draw[arr, dashed] (tree) -- (oracle);
\end{tikzpicture}
\caption{The Pipeline, With the Code Generation This Report Describes}
\label{fig:pipeline}
\end{figure}

\subsection{Refinement of Tasks 1 and 2}
\label{sec:refinement}

Task~1 planned to write the templates against hand-written resolved models, so
that the code generation could proceed while the transformation was still being
written. The first template indeed ran on the hand-written model of the
\texttt{minimal} example. By the time the templates were widened, the Task~2
transformation already resolved every example, so the templates were widened
directly on its output, and the hand-written model is now only the expected
result of the transformation for \texttt{minimal}. No test feeds a hand-written
model to the templates any more.

Generating the files changed the earlier tasks in two ways:

\begin{itemize}
  \item \textbf{The target metamodel gained the files it did not describe.}
    Task~1's metamodel described the workload, identity, network policy,
    route, monitor and secret files. The oracles also hold services, volume
    claims, configuration files, backup jobs, the secret operator's connection
    and the secret store's policies and roles. Seven classes were added for
    them, one class was restructured, and a few values the files need were
    added, such as the endpoint of the release gate and the rule that admitted
    a network peer. The Task~2 report lists every change in its appendix.
  \item \textbf{The application revision became the production
    implementation's.} Every generated Canary carries the revision of its
    application, a digest. Task~2 first computed it over the target model's own
    form, which differs from the production implementation's. Byte-for-byte
    files need the production value, so the transformation now writes the
    application's element in the production implementation's JSON form and
    hashes that.
\end{itemize}

\subsection{Assumptions and Conditions}
\label{sec:assumptions}

\begin{itemize}
  \item The resolved model is complete. The transformation stops rather than
    write a model with a missing value (Task~2), so a template reads every
    value it needs without checking for its absence, except where the target
    metamodel declares a value optional.
  \item The production implementation is the reference for the spelling of
    every file: its key order, its indentation, and when it quotes a string.
    Where YAML allows two spellings, the templates use the production one.
  \item Several resolved models are rendered together when their files share a
    directory. The routes of every project share the directory of their entry
    point, so the index file of that directory lists the routes of every
    project rendered with it.
  \item The templates run without Eclipse, through Acceleo's standalone API, so
    that the build server runs them on every change. A launch configuration
    runs the same module in Eclipse (Appendix~\ref{app:handin}).
\end{itemize}

\subsection{Scope}
\label{sec:scope}

The templates generate every tree the production implementation generates:
the share of the \texttt{minimal} project (9 files), the share of the
\texttt{data} project (28 files), and the files the two share at the level of
the cluster (7 files), 44 files in total. Together they use every kind of file
the specification defines for a running application, apart from two.

The \texttt{auth} and \texttt{knowledge} examples also carry rendered trees, but
those are written by hand as the goal state of the specification: the
production implementation does not generate them yet either, so they hold no
reliable oracle. They contain the two kinds of file no implementation generates
yet, a disruption budget and a database migration job, and the target
metamodel does not describe those two kinds. Section~\ref{sec:limitations}
discusses this.
